Este trabalho tem como objetivo, a partir de um sistema físico chamado
\gls{tecolab}, identificar a dinâmica de uma planta térmica e aplicar um
controlador \gls{lti} em tempo discreto com um tempo de amostragem de
\SI{15}{\second} que seja capaz de seguir uma curva de referência com o menor
erro possível tendo como modelo ideal erro zero em regime permanente.

Partindo desse pressuposto utilizaremos técnicas de identificação da dinâmica da
planta nos sensores, como o ensaio ao degrau de Ziegler-Nichols, e conceitos de
controladores \gls{lti} como \gls{pid} para encontrar uma aproximação com o
menor erro possível a partir de uma curva de referência.

A validação dos métodos aplicados será feita a partir de experimentações
simuladas e ensaio físico no próprio equipamento do \gls{tecolab}.
